% =============================================================================
\chapter{Estado Atual da Implementação}
% =============================================================================

\section{Componentes disponíveis}

Até agosto de 2026, a arquitetura completa ainda não está implementada.
O baseline disponível já reúne a representação simbólica do estado,
pré-condições, efeitos, tarefas primitivas e compostas, métodos de
decomposição, domínio HTN e um \textit{planner} recursivo com \textit{backtracking}.
Ele também inclui o agente executor orientado a \textit{ticks}, sensores para
atualização do estado, replanejamento inicial, uma abstração de \textit{pathfinding}
com BFS para o GridWorld, visualizador em terminal e documentação técnica
inicial.
A seleção de métodos aplicáveis agora é uma extensão explícita por
meio de \texttt{MethodSelectionStrategy}: DFS preserva a ordem declarada do
domínio, a estratégia heurística permite ranqueamento definido pela aplicação,
e a estratégia RL fornece um gancho para uma implementação concreta.

\section{Interpretação correta do estágio atual}

O código existente demonstra a infraestrutura simbólica e operacional
necessária para decompor e executar tarefas.
O baseline executável ainda usa DFS nos exemplos e mantém o backtracking do
planner após a ordenação dos métodos.
O agente conserva uma cópia das tarefas-raiz e a entrega ao planner em cada
tentativa de replanejamento.
Não existe, neste estágio, política MORL treinada, vetor dinâmico de
preferências integrado ao planner ou máscara de métodos utilizada por um
modelo aprendido.

Portanto, o baseline deve ser reportado como infraestrutura inicial, e não como
validação da hipótese completa.

\section{Lacunas para a nova formulação}

Para integrar MORL, a estratégia de seleção deverá transformar os métodos
aplicáveis e o estado simbólico em entradas estáveis para a política, mantendo
uma correspondência inequívoca entre as ações da política e os métodos candidatos.
A integração também exige uma máscara de validade, um vetor de preferências
$\mathbf{w}_t$, objetivos e recompensas vetoriais, além da definição do horizonte
de crédito de cada decisão hierárquica.

Perfis fixos, regras explícitas e codificadores relacionais serão avaliados como
alternativas de representação, e não como requisitos da interface do
\textit{planner}. Logs de decisão e testes unitários e de integração darão suporte
à confiabilidade da implementação. Experimentos automatizados, cenários e sementes
repetidas, baselines de ordem fixa, prioridade estática e utilidade manual, além de
um protocolo estatístico, comporão a validação científica da proposta.

\section{Compatibilidade do baseline com a proposta}

A implementação atual não precisa ser descartada.
Tarefas compostas, métodos, pré-condições, efeitos, sensores e replanejamento
são diretamente reutilizáveis.
A principal mudança arquitetural já realizada é transformar a seleção de método
em um ponto de extensão explícito.
A futura entrada de preferência deverá ser integrada a uma estratégia concreta,
sem substituir as pré-condições simbólicas nem o backtracking do HTN.
A execução de tarefas primitivas continua separada e permanece determinística
no escopo inicial.
